\subsection{The displaced normal cycle}
At infinity a normal to $P$ in sheaf coordinates is $-p$ in the radian coordinates of Section~\ref{sec:coherenttransport}. For $p$ in a cone $\sigma$, the angle on the corresponding face lies in $T_\sigma$ since the facet constants of $P$ are integral. The normal relation on the cover has a canonical parametrization by a ball with an end, using the nearest point of $P$ as in the proof of \eqref{eq:cycledual}. We keep its multiplicities and orientation before projecting to the torus.

\begin{lemma}\label{lem:positivepush}
The class $T_H J_a^{-1}[C_P]$ in the endpoint core-complement group is the small positive Reeb push of the signed polytope normal cycle. The choices work simultaneously for any finite list of such polytopes.
\end{lemma}

\begin{proof}
Along the positive coefficient path from $a$ to $*$, $[C_P]$ is parallel in rightward relative homology. The estimate preceding \eqref{eq:cycledual} applies with all coefficients positive and bounded above and below: every maximizing ray is in the minimal cone of $d\phi_\infty$, its phase is zero on this end, and all uncontrolled phases have a uniform exponential gap. A single truncation puts the entire moving boundary in a common right half-plane. Together with \eqref{eq:relativejunction}, it remains to compute at $*$.

On a normal end write $u=\rho x$, with $\max_v\langle v,x\rangle=1$, $a_v=\langle v,x\rangle$ and $b=\sqrt{2\phi_\infty(x)}$. Displace the angle by $ex/b$ for small $e>0$, cutting off this displacement inside a fixed compact set. The normalized directions form a compact set and $b$ has positive upper and lower bounds there. We need to know on which side of the stopped page this displacement lies, uniformly during the passage to the tailored function. Put $M=e^{\rho-\beta}$ and $m_v=\exp(-\beta+\rho a_v+i\langle v,\theta+ex/b\rangle)$. The formula of \cite[M.6]{FHSZ} is
\[
             f_s=s\eta(1-\chi_0)+
                          \sum_v(1-s\eta(1-\chi_v))m_v .
\]
The weights are real in $[0,1]$. Put $L=\log(2(|V|+1))$, where $V$ is the ray set. On $\rho\ge\beta-L$ a maximizing term has $\chi_v=1$ and $\eta=1$. On every closed normal cone all rays outside the cone have $a_v\le1-\delta$ for a uniform $\delta>0$ by the whole-facet gradient property. For large $\beta$ their cutoffs $\chi_v$ vanish in this band by \cite[M.2 and T.5]{FHSZ}. Their total modulus is at most $|V|(1-s)e^{-\delta\rho}M$. The winning term has angle $e/b$. Every on-cone term has positive real part for uniformly small $e$, and the on-cone terms with negative sine have $a_v<0$, hence total imaginary part of size $O(e e^{-\rho}M)$. The constant term is real and at most $e^LM$. We obtain, with uniform positive constants,
\begin{align*}
 cM&\le\operatorname{Re} f_s\le |f_s|\le CM,\\
 ce-C(1-s)e^{-\delta\rho}&\le \arg f_s
                         \le Ce+C(1-s)e^{-\delta\rho}.
\end{align*}
Also $|\partial_\rho f_0-f_0|\le |V|M/(\mathrm e\rho)$ by maximizing $(1-a_v)e^{-\rho(1-a_v)}$. The first derivative bound for $f_s-f_0$ is $O(e^{-\sqrt\beta}M)$ by \cite[M.8]{FHSZ}, which gives $\partial_\rho\log |f_s|=1+O(1/\rho+e^{-\sqrt\beta})$. On the preceding part $\rho\le\beta-L$, the constant cutoff is one and the monomial sum has modulus below one. Consequently each ray meets $|f_s|=1$ just once, transversely, at $\rho=\beta+O(1)$. Its cut radius is continuous also across the faces of the parametrized link.

First fix a small positive $e_0$, choose $D_->\max b$ with a strict margin as well as the truncation bounds already required, and take $D$ and then $\beta$ large as in Lemma~\ref{lem:equivariantlink}. Enlarge $R=t e^\beta$ for all critical values on both coefficient paths. Now every cut at $e_0$ is below $A_-$, on a literal vertical face of the rounded carrier. Its base angle is positive and small for every $s$. The oriented chains cut on their normal ends thus give a continuous relative transport in the ribbon complement. At $s=0$ cap their boundaries in this positive interval to the page $P_\delta$ of \eqref{eq:relativejunction}, by horizontal transport on the outer circle and, if necessary, along its outer fibre collar. Take the horizontal level large enough to contain this compact cap. It lies in the common right half-plane for $q$ a sufficiently small fixed positive multiple of $R$. In this half-plane $e$ can be decreased to zero on the truncation and the unshifted tail attached there. This proves that the incoming class is the signed $J_*^{-1}[C_P]$. At $s=1$ the lower argument bound is $ce$, so we may reduce $e$ further at this fixed $\beta$.

At $e=0$ the cut in the tailored hypersurface belongs to $H\cap\Gamma$. The Liouville field for each $0\le k\le K$ is transverse there to the source section $|f_1|=1$ and reaches the outer section $\Phi=A$ transversely \cite[U.6--U.9, U.26 and U.31]{FHSZ}. Its first-hit map extends across the core to a local diffeomorphism between the two sections. It is injective in a common neighbourhood of the compact family of cores: otherwise colliding pairs approaching it would have distinct limiting points contradicting core injectivity, or the same limit contradicting local invertibility. Decrease $e>0$ until every cut lies in this neighbourhood. Its base angle makes it disjoint from the core. Attach the outward flow strips to the other section. At $K$ these strips implement exactly the strict change of carrier of Lemma~\ref{lem:equivariantlink}; for all $k$ their outer boundaries miss the fixed outer core. They also miss one sufficiently shrunk full ribbon throughout. Indeed the images of the finite negative intrinsic contractions form nested compact sets in the parameter-boundary product with intersection the family of cores. Our compact family of outer boundaries is disjoint from that intersection. The strips give relative transport through the whole-ribbon deformation.

At $k=0$ the Liouville field in the homogeneous shell is $u\partial_u$ at fixed angle. The final displacement is $e\,dH(p)$ with $H(p)=\sqrt{2\phi_\infty^*(p)}$ by Legendre duality and homogeneity. To pass to a Euclidean unit section, deform $H$ to the dual norm by their convex interpolation $H_l$. For nonzero $p$ we have $\langle p,dH_l(p)\rangle=H_l(p)>0$. At least one generator of the smallest cone containing $p$ therefore has nonzero positive pairing with the displacement; take $e$ small enough that no ray increment reaches a full angular period. The displaced point misses all the pieces of the stop over this momentum. Thus the relative class on the Euclidean section is its positive Reeb push. Projection of each prism from the normal parameter space preserves all multiplicities, also in dimension one.
\end{proof}

\subsection{A cylindrical disk over the dual vector space}
In integral angle coordinates use the reflected momentum $y$ introduced above, a Euclidean norm to identify dual spaces, and primitive $-y\,dm$. A positive push at infinity moves $m$ to $m-e\nu$, $\nu=y/|y|$. Let $P=\{l_i(m)\ge0\}$ with one defining affine function for each facet, and put
\[
      b_P(m)=c-\sum_i\log l_i(m)\ge1,\qquad
      g(y)=\sup_{m\in P^\circ}\{\langle m,y\rangle-b_P(m)\}.
\]
There is a unique maximizer $m=dg(y)$ for every $y$, since the barrier diverges at the boundary and its Hessian is positive definite. The inverse is $y=db_P(m)$. This defines an exact embedded disk, a graph over momentum even after its base is projected to the torus. We describe a cylindrical representative $B_P$ of its positive completion and determine both its geometric class and its class under microlocalization.

For $y=r\nu$ and $m=dg(y)$, every $z\in P$ satisfies
\[
             \langle y,z-m\rangle=\sum_i(1-l_i(z)/l_i(m))
                              \le \#\{\text{facets}\}.
\]
Testing at a fixed interior point also bounds $1/l_i(m)$ by $C(1+r)$ for all $i$, so $b_P(m)=O(1+\log(1+r))$. If $h_P$ denotes the support function, $g(r\nu)/r$ converges uniformly to $h_P(\nu)$ and every limiting $m$ lies on the face exposed by $\nu$. Let $\chi$ be a smooth nondecreasing cutoff, zero before one and one after two. At a large radius $B$ set
\[
        g_B(r\nu)=(1-\chi(r/B))g(r\nu)
                         +\chi(r/B)\,r g(B\nu)/B .
\]
Subtract $e r$ with an inner cutoff, zero near the origin and one well before this large annulus. Its differential is the graph of $B_P$. It is exactly conical past $2B$ by homogeneity. On $B\le r\le2B$, the extra radial derivative from $\chi$ is $o(1)$ by the uniform support-function limit. Replacing $dg(B\nu)$ by the derivative of $r g(B\nu)/B$ changes its radial component by $-b_P(dg(B\nu))/B=o(1)$. Thus all accumulation points of the undisplaced ends are in the exposed face, by convexity also in the interpolation region. For small fixed $e>0$ their displaced limiting links miss the fan stop by the last argument of Lemma~\ref{lem:positivepush}. Take $B$ large enough for actual disjointness.

There is a corresponding relative homology from this disk to the normal cycle. In the angle cover parametrize by $w=m+y^\sharp$. For $0\le d\le1$ minimize $|w-m|^2/2+d\,b_P(m)$ over $P$, using nearest-point projection when $d=0$. The minimizer is continuous down to zero by testing the minimum against each fixed interior point and using $b_P\ge1$. The preceding inequality gives $\langle y,z-m\rangle\le d\,\#\{\text{facets}\}$ on all these graphs. On a large $w$-sphere all ends are therefore uniformly close to the normal relation. Displace them by $-e\nu$ and clamp large momenta radially to a common boundary section. The prism on a single large $w$-ball has side in the stop complement. The same is true while conifying the smooth terminal graph by the interpolation above, now parametrizing by momentum. At $d=1$ the coordinate change preserves orientation since $I+D^2g$ is positive. This proves that the relative class of $B_P$ is the displaced signed normal cycle up to choice of orientation. In the ordinary positive ray torus the same prisms can instead be truncated in the right half-plane. Near each point of the compact normal link, all exponents close to the maximum have phases close to zero, and the others have a uniform strict gap. The real part tends uniformly to infinity on these neighbourhoods as above. A sufficiently small displacement can therefore be removed in the rightward group. Both comparisons use the same orientation. We determine it next.
